
\subsubsection{3.1 Error Categorization}

Following RLTF's categorization, we distinguish:

\textbf{U\_line errors:} SyntaxError, IndentationError, NameError, TypeError, AttributeError, KeyError, IndexError, ValueError, ZeroDivisionError. These exceptions report the exact line where the error occurs.

\textbf{U\_ignore errors:} RuntimeError, RecursionError, MemoryError, TimeoutError, AssertionError. These exceptions often report symptom lines rather than root causes.

\subsubsection{3.2 Gating Mechanism}

Standard RLTF applies fine-grained penalties to all errors with traceback information. For a generated code sequence with error at line $\ell$, tokens at line $\ell$ receive penalty -1.0 while other tokens receive -0.1:

\begin{verbatim}
r_fine(token_i) = -1.0 if line(token_i) == ℓ else -0.1
\end{verbatim}

The gated approach conditions this on error type:

\begin{verbatim}
r_gated(token_i, error_type) = 
    r_fine(token_i)   if error_type ∈ U_line
    r_coarse          if error_type ∈ U_ignore
\end{verbatim}

For U\_line errors, the standard fine-grained penalty is applied. For U\_ignore errors, only the coarse reward signal is applied.

\subsubsection{3.3 Implementation}

The modification to RLTF reward computation:

\begin{enumerate}
\item Execute generated code in sandbox
\item If error occurs, parse traceback for exception type and line number
\item Classify exception type as U\_line or U\_ignore
\item Apply fine-grained penalty only if U\_line; otherwise coarse-only
\end{enumerate}

The categorization lookup is O(1)---a set membership check.

